
\subsubsection{3.1 Task Selection}

We selected dual-sensitive programming tasks from the HumanEval benchmark augmented with HumanEval+ extended test suites. A task qualifies as dual-sensitive if, across K=20 code samples generated by CodeLlama-7B:

\begin{itemize}
\item At least one sample fails static analysis (mypy --strict) but passes execution tests
\item At least one sample passes static analysis but fails execution tests
\item Within-task variance SD $\leq$ 1.0
\end{itemize}

Out of 164 HumanEval tasks, 35 (21.3\%) met these criteria. These tasks exhibit error patterns detectable by both feedback sources, enabling paired experimental design.

\subsubsection{3.2 Experimental Conditions}

\textbf{CASCADE (Sequential Routing):}
\begin{enumerate}
\item Generate code sample
\item Apply mypy --strict static analysis
\item If mypy fails: provide only mypy feedback, return to step 1
\item If mypy passes: run pytest execution tests
\item If pytest fails: provide only pytest feedback, return to step 1
\item If both pass: task solved
\end{enumerate}

\textbf{AGGREGATION (Simultaneous Baseline):}
\begin{enumerate}
\item Generate code sample
\item Apply both mypy --strict and pytest simultaneously
\item Concatenate feedback from both sources
\item If either fails: provide combined feedback, return to step 1
\item If both pass: task solved
\end{enumerate}

Both conditions use identical model configuration (CodeLlama-7B base, temperature=0.8, top-p=0.95, top-k=40, max tokens=256) and are limited to 10 iterations per task.

\subsubsection{3.3 Static Analysis Configuration}

Static type checking uses mypy version $\geq 1.5.0$ with --strict mode, which enforces:
\begin{itemize}
\item Explicit type annotations for function signatures
\item Strict optional checking (None handling)
\item No implicit Any types
\item Disallowing untyped function definitions
\end{itemize}

Timeout: 10 seconds per verification.

\subsubsection{3.4 Execution Testing Configuration}

Runtime testing uses pytest version $\geq 7.0.0$ with HumanEval+ augmented test suites (80+ tests per task). Tests are executed in isolated temporary files with 120-second timeout per sample.

\subsubsection{3.5 Metrics}

\textbf{Primary metrics:}
\begin{itemize}
\item Mypy error detection rate: proportion of iterations where mypy identifies errors before execution
\item Token efficiency ratio: (CASCADE tokens per task) / (AGGREGATION tokens per task)
\end{itemize}

\textbf{Secondary metrics:}
\begin{itemize}
\item Gating efficiency: percentage of iterations where pytest execution is skipped in CASCADE
\item Token breakdown: mypy vs. pytest contribution for both conditions
\item Success rate: proportion of tasks solved within iteration limit
\end{itemize}

\subsubsection{3.6 Hypothesis Structure}

The investigation tested four sub-hypotheses:

\textbf{H-E1 (EXISTENCE):} Dual-sensitive tasks exist in sufficient quantity (N $\geq$ 20) with adequate variance (SD $\leq$ 1.0) for experimental design.

\textbf{H-M1 (MECHANISM):} Static analysis detects $\geq 30$\% of errors instantly with zero execution cost.

\textbf{H-M2 (MECHANISM):} Sequential single-source feedback reduces iteration count via attention economy.

\textbf{H-M3 (MECHANISM):} Conditional execution gating maintains token efficiency within 15\% of aggregation baseline.

Gate criteria: H-E1 and H-M1 classified as MUST\_WORK (failure blocks downstream hypotheses). H-M2 and H-M3 classified as SHOULD\_WORK (failure documented as limitation).
